\begin{enumerate}[resume]
    \item In this chapter you have to present all of the details of the short paragraphs you spoke about to test your research question in the methodology section.   
    %Clarity:
    \item  Your results needs to be clear and well organised.
    %Result and problem statement:
    \item  Demonstrate how the collected results answers the research questions. 
    %Results and Discussion:
    \item  Consider breaking your results section into two parts. First try to look at your graphs and write down that you see (observations). Thereafter take all your observations and try to infer bigger questions or arguments in the next chapter (inference/discussion). You can put these larger inference in the discussion section.
    %Realised limitations from mnethodology
    \item  Please include a section where you discuss how the limitations that you state on the methodology sections impacts the obtained results. 
    %More Graphics:
    \item  Can you include more graphics to describe your result? Remember a picture can say 1000 words (although you will still have to describe what you feel in important in the graphic).
    %Severity of Misclassifications:
    \item  If applicable, can you account for the `type' of misclassifications which occurred? Do you think that some misclassification are more severe than others?    
    
    
    %Hypothesized outcome vs actual outcome:
    \item  What kinds of results were expected in comparison to what was achieved? (tool, framework, new theory, new technical solutions, etc).
    %Applications:
    \item  Practical relevance of the expected results (How to apply the results?).
    %More depth:
    \item  The results sections needs to include the findings of the study. These include data presented in tables, charts, graphs, plots, and other analysis type figures. There should also be a contextual explanation of all the figures. \href{https://blog.wordvice.com/writing-the-results-section-for-a-research-paper/#:~:text=The\%20Results\%20section\%20should\%20include,its\%20meaning\%20in\%20sentence\%20form}{Here is a resource for writing results}.
    %Neaten-up Results: 
    \item Try to neaten-up the results by perhaps incorporating them into figures or tables? Expressing the results better can save space and make the work easier to read.
    %Confidence Intervals and Error Bars:
    \item  You need to incorporate some confidence intervals in your results section. To do this you need to run multiple trials of the predictive model (like 10 trials), thereafter you can plot the mean of these trials and use the standard deviation as error bars. You can also tryout other methods like rolling means with standard errors, but there should be some explicit indication of uncertainty.
\end{enumerate}